% intro.tex
\clearpage
\begin{center}
The word concerning\\
Samuel, Saul and David

Comprising the books commonly referred to in English Bibles as\\
First Samuel, Second Samuel and part of First Kings
\end{center}
{ % keep parindent values local
\setlength{\parindent}{1.5em}
\section*{Foreword}
\subsection*{The Reason for this Edition}

There are a number of copies of the various books that make up what is called in English the Bible.
Some of these copies are faithful but some,
sadly, are less so. Differences between copies can make it difficult to know
which truly represent the Word of God. There are tests, however, which
we can use to recognise the true word of God. These are the tests that
God has given for his people to use to discern if a person who claims to
speak in His name does so truly or not.

For the Hebrew scriptures (the Law, the Prophets and the Writings), the 'standard' copies go by the name the ``Masoretic'' Text, or MT for short. These are the copies that virtually all English translations are based on, and many other translations, too. When Jerome translated the Latin Vulgate in the late 4th century he did so from the MT copies that had been provided to him by Jewish scholars. For most of the Hebrew scriptures, these copies can be shown to be substantially faithful and correct. However, that is not the case with the Word concerning
Samuel and Kings Saul and David comprising the books that go by the names of First and Second Samuel and a little of First Kings. Here it can be shown that these
`standard' copies again and again fail the test of faithfulness. In some places the standard copies are virtually incomprehensible and in others they contradict other teachings of scriptures.
Fortunately, this family of copies is not all we have. There exists a
family of Hebrew and Greek copies that are sometimes called the
``Palestinian'' copies. These copies include the ``Lucianic'' Greek copies and the
Hebrew Dead Sea Scrolls of Samuel, particularly that called 4Q51, also
called 4QSam\textsuperscript{a}. The table below lists a few of the
differences between the Palestinian and Masoretic copies with evidence
for the reliability of the former over the latter.
\begin{tabular}{p{0.1\textwidth}p{0.25\textwidth}p{0.25\textwidth}p{0.3\textwidth}}\hline
  Ref.& Palestinian text& Masoretic Text&Note\\\hline
  1 Sam 13:1&he reigned forty years&he reigned two years&See the
                                                          footnote in
                                                          the WEB
                                                          text. Compare
                                                          Acts
                                                          13:21.\\
  1 Sam 13:1&Saul was thirty years old&Saul was years old&The number
                                                           is omitted
                                                           in the MT
                                                           (again, see
                                                           the
                                                           footnote in
                                                           the WEB)\\
  1 Sam 13:8&a donkey loaded with an omer of bread&a donkey of
                                                    bread&The WEB
                                                           interpolates
                                                           where the
                                                           MT makes no
                                                           sense, this
                                                           time with
                                                           no footnote.\\
  1 Sam 18:27&a hundred men&two hundred men&Compare 2 Sam 3:14\\
  2 Sam 8:4&1,000 chariots and 7,000 horsemen&one thousand seven
                                               hundred
                                               horsemen&Compare 1 Chr
                                                         18:4\\
  2 Sam 8:13&18,000 Edomites&eighteen thousand men of the
                              Syrians&Syrians (Arameans) is similar in
                                       Hebrew to Edomites. Compare 1
                                       Chr 18:21\\
  2 Sam 15:7&at the end of four years&At the end of forty
                                       years&Compare 1 Kings 2:11, 1
                                              Chr 3:1--4\\
  2 Sam 21:8&five sons of Merab&five sons of Michal&The WEB renders
                                                     the P text here
                                                     to avoid the
                                                     obvious error in
                                                     the MT. Compare 1
                                                     Sam 18:19, 2 Sam
                                                     6:23.
\end{tabular}
The Palestinian copies are not a really rare or unusual form of
Samuel. It is the version quoted by Josephus and many so-called
ante-Nicene fathers. But it is not easy to obtain in English, which is the reason
for this edition.

Some years ago the editor of this edition undertook a translation into
English of these more faithful copies of this part of scripture. This
edition presents this translation in the left column in parallel with
the World English Bible translation (broadly based on the MT, but with
  some deviation as noted above) in the right
column. This translation is based heavily on the American Standard
Version of 1901, as noted below, adjusted to match the Palestinian
copies. Despite my best efforts, it is likely that there have been
some (hopefully minor)
differences between the P text and the ASV English that have been
overlooked.

It is worth commenting on one particular type of difference between
the families of copies. Where the Palestinian text includes names like
Mephibaal (see 2 Sam 4:4, for example), the MT has Mephibosheth. The
-bosheth suffix means `shame' and this is clearly a scribal euphemism
to avoid naming the Canaanite false god Baal. The corresponding names
in Chronicles retain the -baal form that is seen in the Palestinian
copies.

\subsection*{The Name of This Document}

This edition contains all of what is traditionally called "First Samuel"
and "Second Samuel", and two chapters of what is traditionally called
"First Kings". These traditional names are somewhat misleading, particularly as the death of
Samuel is recorded in the first part (25:1). Even the division into a
'first' and a 'second' book is somewhat arbitrary, as is the
distinction between the books of `Samuel' and the books of
`Kings'. (In the Greek Septuagint translation of the Hebrew
scriptures, all four of the books of Samuel and Kings go by the name
`\textgreek{Βασιλειῶν},' that is, `of Kingdoms.')
I have called this document 
``Samuel, Saul and David Parallel Edition'' as this is descriptive of the
content. 

In previous editons of my translation I tried to avoid using traditional names where they might be considered misleading. For instance, the books of the Law were all referred to as Moses with a Roman numeral to indicate which book was being referenced.  In this edition I have reversed that practice as making those references unnecessarily obscure to readers with less familiarity with the Bible. There are some footnotes that still retain this more particular approach to naming parts of scripture. References in footnotes to ``first part'' or ``second part'' refer to the parts of this document corresponding to 1 Samuel and 2 Samuel.

\subsection*{Verse and Chapter Divisions}

Different editions use slightly different verse numbering systems. The
reader should remember that these divisions are entirely arbitrary. This
edition follows, as much as possible, the numbering system used in the
majority of English language editions. In any case, the parallel
presentation should make any variations clear.

The chapter numbering of the third section follows the traditional
numbering system which has ``First Kings'' commencing at that
point. Some other editions count the text included there as chapters
25 and 26 (of ``2 Samuel'') 


\subsection*{Principles of Translation}

I have favoured a word-for-word approach in preparing my translation,
translating the words and grammatical structures from the original
language to English, but leaving the idioms and linguistic style
untranslated. This jumps out at the reader from the very first word:
'And'---a word not generally used to begin an account! Indeed, the
Hebrew conjunction \texthebrew{ו} could have been reasonably translated with many
different words ('now', 'so', 'but', etc.) or left out entirely where
English usage does not require it, but I have sought to make this not
only readable, but usable for study where details like the actual words
used in the original may be of interest.

Text presented \textit{in italics} represents text not found in the original
language, but needed to make the English text understandable.

\subsection*{'You'---Singular/Plural Ambiguity}

Sometimes it is not possible with modern English to distinguish between
you (singular) and you (plural), even though the original language text
is unambiguous. Wherever I have noticed that the English translation is
ambiguous as to whether the word 'you' refers to one person or more than
one I have resolved the ambiguity for the reader unobtrusively by using
a single dagger for singular and a double dagger for plural, thus:
you\textsuperscript{†} (singular); you\textsuperscript{‡} (plural).

\subsection*{The Name of God}

The name of God is rendered in the Palestinian translation as ``the \textsc{Lord}''. In some places this is
used with the ordinary word ``lord''. Although unambiguous in print,
when reading aloud one may choose to use ``Jehovah'' or
``\textsc{God}'' instead.

In the WEB this name is rendered as `Yaweh', a transliteration of the
Hebrew.

\subsection*{Acknowledgements}

The translation of the reliable copies was inspired by and relies heavily on a manuscript comparison
prepared by Evan Blackmore. Translation is loosely based on word choice
and syntax of the American Standard Version 1901 but extensively updated
to reduce the archaisms characteristic of the ASV, and with extensive
reference to Hebrew and Greek texts.

The World English Bible is published by eBible.org and can be found at
\href{https://ebible.org/find/show.php?id=eng-web}{https://ebible.org/find/show.php?id=eng-web}. It
has been explicitly placed in the public domain.

} % keep parindent values local

\bigskip
\noindent Glen Prideaux\\
March 2026

\bigskip
\noindent\textbf{Left:} Prideaux translation of Palestinian copies,
labelled \circled{\small P} and\\
\textbf{Right:} World English Bible based (mostly) on the Masoretic copies,
labelled {\small (WEB)}